\scanpage{97}

\[
\underbrace{-\frac{r_2}{r+1},\ldots,-\frac{r_2}{r+1}}_{r_1},
\underbrace{\frac{r_1}{r+1},\ldots,\frac{r_1}{r+1}}_{r_2}
\qquad
[1,r+1]=\underbrace{[1,r_1]}_{I_1}\cup\underbrace{[r_1+1,r+1]}_{I_2}
\]
\[-\frac{r_1}{r+1}\quad-\frac{r_2}{r+1}\qquad\pm1\qquad \sigma^{\unsure{-1}}I_1,\ \sigma I_2\]

\[
\begin{aligned}
\xi&\qquad\alpha_1\text{ sur }I_1,\ \alpha_2\text{ sur }I_2\\
\sigma.\xi&=\xi\circ\sigma^{-1}\qquad\alpha_1\text{ sur }\sigma I_1,\alpha_2\text{ sur }\sigma I_2\\
\xi-\sigma\xi&\qquad 0\text{ sur }(I_1\cap\sigma I_1)\cup(I_2\cap\sigma I_2)\\
&\qquad-1\text{ sur }I_1\cap\sigma I_2\\
&\qquad 1\text{ sur }I_2\cap\sigma I_1
\end{aligned}
\]

On veut, $\forall\sigma$ tel que $\sigma I_1\ne I_1$ \inserted{et $\sigma I_2\ne I_2$}, \deleted{\unsure{que}} \inserted{i.e. $\sigma I_2\cap I_1\ne\varnothing$} i.e. $\sigma I_1\cap I_2\ne\varnothing$

\deleted{$I_1\cap\sigma I_2$}

\[
\begin{aligned}
\operatorname{card}(I_1\cap\sigma I_2)&=1\ \to\ \text{c’est possible \inserted{(pour tout $\sigma$)} ssi }\operatorname{card}I_1\leq1\text{ ou card }I_2\leq1\\
\operatorname{card}(I_2\cap\sigma I_1)&=1\qquad\rule{1em}{.3pt}\quad\rule{1em}{.3pt}
\end{aligned}
\]

\deleted{Donc} OK ssi $r_1=1$ ou $r_2=1$

\[
\begin{aligned}
&-\frac1{r+1},\underbrace{\frac r{r+1},\ldots,\frac r{r+1}}_r\\
&\underbrace{-\frac r{r+1},\deleted{\unsure{\tfrac1{r+1}}}\ldots,-\frac r{r+1}}_r,\frac1{r+1}
\end{aligned}
\]

2 orbites \inserted{(si $r\geq2$)} \unsure{pour|sous} $W$, qui sont \deleted{\unsure{conjuguées par}} \inserted{transformées} l’une de l’autre par \unsure{l’élém|l’élément} \deleted{\unsure{non triv}} de $\unsure{\Gamma|E}=\widehat W/W$ i.e. par la symétrie !

$\text{\textcircled{$B_r$}}\qquad V=E=\mathbb R^r$

\[
\begin{aligned}
R&=\{\pm\varepsilon_i\ (1\leq i\leq r)\quad\pm\varepsilon_i\pm\varepsilon_j\ (1\leq i<j\leq r)\}\\
&2r+4\frac{r(r-1)}2=2r+2r(r-1)=2r^2
\end{aligned}
\]

\[
2r+1\quad
\begin{pmatrix}\lambda&&a_{ij}\\-a_{ij}&\ddots&\\&&0\end{pmatrix}
\qquad
\begin{pmatrix}\lambda_1\\&\lambda_1^{-1}\\&&\ddots\\&&&\lambda_r\\&&&&\lambda_r^{-1}\\&&&&&1\end{pmatrix}
\]

Produit scalaire habituel, longueurs 1 et 2

\[
\begin{aligned}
\xi&=(\xi_1,\xi_2\ldots\xi_r)\quad\unsure{\forall\xi|\deleted{\xi_i}}\quad0\leq\xi_1\leq\xi_2\leq\ldots\leq\xi_r\\
&\underbrace{\xi_1\ldots\xi_{r_1}}_{\alpha_1}\quad\underbrace{\xi_{r_1+1}\ldots\xi_{r_1+r_2}}_{\alpha_2}\ \ldots\
\underbrace{\xi_{r_1+\cdots+r_{s-1}+1}\ldots\xi_{r_1+\cdots+r_s=r}}_{\alpha_s}\\
(\xi,\varepsilon_i)&\in 2\cdot\{-1,0,1\}\qquad\text{i.e. }\boxed{\xi_i\text{ égale }0\text{ ou }2}\\
(\xi,\varepsilon_i\pm\varepsilon_j)&=\xi_i\pm\xi_j\in\{1,0,-1\}\\
&\deleted{\unsure{\xi}}\quad\alpha_i+\alpha_j=1\quad\deleted{\unsure{\text{i.e.}}}\\
&\phantom{x}\quad\alpha_i-\alpha_j=-1\qquad\text{si }1\leq i<j\leq s\\
&\hspace{12em}\unsure{\forall}
\end{aligned}
\]

\deleted{\unsure{bien seule|Une seule} donc} impossible, donc $s=1$

\[\xi=(2,2,\ldots,2)\qquad\unsure{\mathscr S|\mathscr G|\mathfrak S}=\varnothing\]

$\text{\textcircled{$C_{\unsure{r|6}}$}}\qquad E=\mathbb R^r$

\[
R=\{\pm2\varepsilon_i\ \inserted{(1\leq i\leq r)},\ \pm\varepsilon_i+\varepsilon_j\}\ (1\leq i\leq j\leq r))
\]

\deleted{\unsure{ce}}\qquad $\deleted{\unsure{\xi_i\in}}$

\[
\begin{aligned}
(\xi_i)_{1\leq i\leq r}&\qquad \xi_i\in\tfrac24\{-1,0,+1\}=\{\\
0\leq\xi_1&\leq\ldots\leq\xi_r\qquad\text{i.e. }\xi_1\in\{0,\tfrac12\}\\
\pm\xi_i\pm\xi_j&\in\{-1,0,+1\}\\
\deleted{\unsure{\theta}}&(\tfrac12,\tfrac12\ldots,\tfrac12)\qquad\deleted{r\geq3}\quad r\geq3\\
\mathscr S&=\deleted{\unsure{\tfrac12\{\varepsilon_1,\varepsilon_2,\ldots\}}}\ \deleted{\unsure{\tfrac12}}\{\pm\tfrac12\}^r\qquad\unsure{=}\\
\mathscr S_0&=\varnothing
\end{aligned}
\]

$\text{\textcircled{$D_r$}}\qquad E=\mathbb R^r$

\[
\begin{aligned}
R&=\{\pm\varepsilon_i\pm\varepsilon_j\quad1\leq i<j\leq r\}\\
&\deleted{\unsure{\xi}}\quad \xi_1\ldots\xi_r\\
&\xi_1\leq\underset{\geq0}{\xi_2}\ldots\leq\xi_r\quad \pm\xi_i\pm\xi_j\in\{-1,0,1\}
\end{aligned}
\]

au plus \deleted{deux} valeurs distinctes des $\xi_i$, différant de 1, soit $\alpha$ et $\alpha+1$, \inserted{mais \deleted{\unsure{on|$\alpha+$}} $2\alpha+1\in\{-1,0,1\}$ i.e.} ou bien ttes les valeurs égales, égales : $\pm\frac12$

\[
\begin{aligned}
\deleted{2\alpha+1=0}&\quad\alpha=-\tfrac12\quad\alpha+1=\tfrac12
&-\tfrac12,\underbrace{\tfrac12,\tfrac12\ldots\tfrac12}_{r-1}\\
2\alpha+1=1&,\ \alpha=0&(\underbrace{0,\ldots0}_{\unsure{r-1|r+1}},\underbrace{1}_{r_2=1})\ \deleted{\unsure{1-}}\ \to r\\
2\alpha+1=-1&,\ \alpha=-1&(-1,\underbrace{0,\ldots,0}_{r-1})
\end{aligned}
\]
\deleted{\unsure{au facteur}}\quad à conj. près $(r_2=1)$

\[
\begin{aligned}
&\bigl(\tfrac12(\unsure{\varepsilon|\pm})_i\bigr)\qquad
\deleted{\unsure{\{\pm1\}}}\quad\prod(\unsure{\varepsilon})_i=1
\ \Big\}\cdot(\tfrac12,\tfrac12,\ldots,\tfrac12)\\
&\ldots\qquad\qquad\rule{1em}{.3pt}\qquad\prod(\unsure{\varepsilon})_i=-1\\
&\deleted{\unsure{\varepsilon_i}}\{\pm\varepsilon_i\}
\end{aligned}
\]

\deleted{il y a \unsure{trois}} $\mathscr S$ a trois orbites\\ $\mathscr S_0=\varnothing$ \deleted{\unsure{et a une orbite}}

% archive sheet number: 98
